\section{Fundamentos de Mechanistic Interpretability}

\subsection{¿Qué es Mech Interp?}

Las redes neuronales «crecen en lugar de programarse»: la arquitectura es el andamiaje, la función de pérdida es la luz guía, y el descenso de gradiente produce un artefacto que no sabemos cómo programar directamente. \textit{«No las fabricamos, las cultivamos... es casi una entidad biológica o un organismo que estamos estudiando.»}

Aquí hay dos preguntas entrelazadas: una pregunta científica profunda (¿qué ocurre en realidad dentro de estos sistemas?) y una pregunta de seguridad crucial (¿cómo garantizar que sean confiables?).

\begin{importantbox}{No es atribución, es ingeniería inversa}
Un Saliency Map (por ejemplo, «qué parte de la imagen hace que el modelo crea que es un perro») \textbf{no es} Mechanistic Interpretability. Lo que busca Mech Interp son \textbf{algoritmos y mecanismos}: aplicar ingeniería inversa a los pesos de la red neuronal para obtener algoritmos comprensibles.

Analogía: los pesos de la red neuronal equivalen a un programa binario de computadora. El objetivo es descompilar este «programa compilado» en un algoritmo legible. Los valores de activación equivalen a la memoria; los pesos, a las instrucciones; ambos deben entenderse.

Actitud central: \textit{«el descenso de gradiente es más inteligente que tú»}: descubrimiento de abajo hacia arriba, no hipótesis de arriba hacia abajo. No presupongas qué hay dentro del modelo, descúbrelo.
\end{importantbox}

\subsection{Universality: consistencia entre redes e incluso entre especies}

Modelos de visión con arquitecturas distintas forman las mismas características: filtros de Gabor, detectores de curvas, detectores de frecuencia alta y baja. Más sorprendente aún: estas características también existen en redes neuronales \textbf{biológicas}:

\begin{knowledgebox}{Consistencia de características entre especies}
\begin{itemize}
    \item \textbf{Detectores de curvas}: descubiertos primero en redes neuronales artificiales, confirmados después en el cerebro de monos
    \item \textbf{Detectores de frecuencia alta y baja}: descubiertos primero en redes neuronales artificiales, confirmados después en el cerebro de ratones
    \item \textbf{Neurona de Donald Trump}: cada modelo de visión tiene una neurona dedicada a Trump, que responde tanto a imágenes de su rostro \textbf{como} a la palabra «Trump». Es un concepto abstracto, no solo coincidencia de patrones
\end{itemize}

\textit{«El descenso de gradiente, en cierto sentido, encuentra la forma correcta de dividir las cosas... muchos sistemas convergen hacia las mismas abstracciones.»} Esto sugiere que existe una especie de «frontera natural de abstracción» que se descubre sin importar si el sustrato de cómputo es silicio o carbono.
\end{knowledgebox}

\subsection{Resumen de la sección}

Mechanistic Interpretability busca entender el «algoritmo» de la red neuronal, no solo atribuirlo. La consistencia de características entre redes y entre especies sugiere que el descenso de gradiente está descubriendo una especie de «forma natural de dividir las abstracciones».
